\documentclass[border=4pt]{standalone}
\usepackage{amsmath,amssymb,mathtools,xcolor,varwidth}
\definecolor{tangreen}{HTML}{00A35A}
\begin{document}
\begin{varwidth}{28em}
\large\bfseries
Draw the \textcolor{tangreen}{tangent $ZPR$ at $P$ — the straight path the body WOULD take if the force suddenly      vanished}; this line is the heart of the argument. The distance by which the curve pulls away from that      tangent is the versed sine $QR$, drawn from the nearby point $Q$ back toward the tangent. By Cor. 4 of      Prop. I this fall in a given time is directly as the force. And by Cor. 2 and 3 of Lem. II, lengthening the      time in any ratio augments the arc in that ratio and the versed sine in the DUPLICATE ratio — so the versed      sine is as the force times the square of the time.
\end{varwidth}
\end{document}
